% ══════════════════════════════════════════════════════════════════════════════
\chapter{RF Front-End}
\label{ch:rf-frontend}
% ══════════════════════════════════════════════════════════════════════════════
\section{RF PCB Layout Considerations}
\label{sec:rf_layout}

The performance of the RF front-end is strongly influenced by the PCB layout. While both the 144\,MHz (VHF) and 433\,MHz (UHF) bands are relatively low in frequency compared to microwave systems, careful RF design is still required to ensure impedance control, minimize losses, and reduce electromagnetic interference.

On the first RF front-end prototype, a 4-layer stackup was considered sufficient because it had no size constraints. However, to achieve a more compact design within the form factor of a Raspberry Pi shield, a 6-layer stackup was adopted. This enables improved separation of RF, power, and digital domains while maintaining a continuous ground reference.

\subsection{6-Layer Stackup Strategy}

The PCB is implemented using JLCPCB's \texttt{JLC0612IH-3313} 6-layer stackup with a nominal thickness of 1.2\,mm. The layer arrangement is selected to ensure that RF traces on the top layer have a well-defined impedance with a continuous ground reference directly beneath them, while also providing adequate separation between RF, power, and digital domains.

\begin{table}[H]
    \centering
    \caption{6-layer PCB stackup with electrical function and material properties.}
    \label{tab:rf_stackup_functional}
    \small
    \renewcommand{\arraystretch}{1.2}
    \begin{tabularx}{\textwidth}{l l l l c X}
        \toprule
        \textbf{Layer} & \textbf{Type} & \textbf{Thickness} & \textbf{Material} & \textbf{$\varepsilon_r$} & \textbf{Function} \\
        \midrule
        F.Cu   & Copper       & 0.035\,mm (1\,oz)   & ---     & ---  & RF routing, matching networks, SMA feedlines \\
        Dielectric 1 & Prepreg 3313 & 0.0994\,mm & NP-155F & 4.1  &  \\
        In1.Cu & Copper       & 0.0152\,mm (0.5\,oz)& ---     & ---  & Primary ground plane (RF return path) \\
        Dielectric 2 & Core   & 0.35\,mm            & NP-155F & 4.36 &  \\
        In2.Cu & Copper       & 0.0152\,mm (0.5\,oz)& ---     & ---  & Power distribution (3.3\,V, 5\,V) \\
        Dielectric 3 & Prepreg 2116 & 0.1164\,mm & NP-155F & 4.16 &  \\
        In3.Cu & Copper       & 0.0152\,mm (0.5\,oz)& ---     & ---  & Secondary ground plane (coupling reduction) \\
        Dielectric 4 & Core   & 0.35\,mm            & NP-155F & 4.36 &  \\
        In4.Cu & Copper       & 0.0152\,mm (0.5\,oz)& ---     & ---  & Digital and control signal routing \\
        Dielectric 5 & Prepreg 3313 & 0.0994\,mm & NP-155F & 4.1  &  \\
        B.Cu   & Copper       & 0.035\,mm (1\,oz)   & ---     & ---  & Auxiliary signals and component placement \\
        \bottomrule
    \end{tabularx}
\end{table}

This stackup ensures that RF traces on the top layer (F.Cu) are implemented as controlled-impedance microstrip lines referenced to the continuous ground plane on In1.Cu. The relatively thin dielectric (0.0994\,mm) enables a compact 50\,$\Omega$ trace width (0.15\,mm), while the additional ground plane (In3.Cu) improves electromagnetic shielding and reduces coupling between the RF, power, and digital domains.

The use of dual ground planes (L2 and L4) improves electromagnetic shielding and reduces coupling between RF and digital layers, while separating digital routing on L5 minimizes interference with sensitive RF signals on the top layer.

The key rule is that all RF traces on L1 must have an uninterrupted ground reference directly beneath them on L2. No splits or voids are allowed under RF lines, as this would disrupt the current return path, increase loop area, and lead to additional radiation and impedance discontinuities.

\subsection{Transmission Line Implementation}

All RF signal paths were routed with controlled impedance considerations, targeting an impedance of 50\,$\Omega$.

In theory, grounded coplanar waveguide (GCPW) is preferred for RF routing due to improved field confinement, reduced radiation, and better isolation. This was implemented on the first prototype. However, due to minimum via size constraints (to reduce PCB manufacturing cost), the matching network occupied a significant area.

In the final design, microstrip routing was used for both the 144\,MHz and 433\,MHz bands. This choice was primarily driven by space constraints imposed by the compact form factor (Raspberry Pi shield size), as GCPW requires additional lateral clearance for ground pours and via fences, increasing the required routing area.

At the given frequencies, where wavelengths are approximately 2.08\,m and 0.69\,m respectively, and considering that the complete matching network fits within approximately 10\,mm $\times$ 15\,mm, these dimensions are small compared to the wavelength. In particular, they remain well below \(\lambda/10\) (approximately 208\,mm and 69\,mm at 144\,MHz and 433\,MHz, respectively), meaning the interconnects can be considered electrically short. As a result, the expected performance difference between GCPW and microstrip is small, making microstrip a practical and sufficiently accurate solution while allowing a more compact layout.

\subsection{Frequency-Dependent Layout Considerations}

Although both frequency bands share the same layout principles, their sensitivity to PCB effects differs.

\subsubsection{144\,MHz (VHF)}

At 144\,MHz, the wavelength is relatively large, and PCB traces exhibit more lumped behavior.

\begin{itemize}
    \item Transmission line effects are moderate but still considered for all RF paths.
    \item Microstrip routing is generally sufficient for most connections.
    \item Layout is relatively tolerant to small discontinuities.
    \item Via stitching density can be moderate (e.g., 3--5\,mm spacing).
\end{itemize}

\subsubsection{433\,MHz (UHF)}

At 433\,MHz, the wavelength is shorter, and PCB parasitics significantly affect performance.

\begin{itemize}
    \item All RF traces must be treated as transmission lines.
    \item GCPW is preferred for improved impedance control and isolation.
    \item RF traces must be kept short, straight, and continuous.
    \item Via stitching should be dense (e.g., 1--2\,mm spacing) to maintain a strong ground reference.
    \item Layout is highly sensitive to discontinuities, stubs, and parasitic effects.
\end{itemize}

Although the electrical behavior of the 144\,MHz (VHF) and 433\,MHz (UHF) bands differs, both frequency bands are routed using the same PCB design strategy. The layout is designed according to the more stringent requirements of the UHF band, ensuring that the VHF performance is inherently satisfied.

\subsection{Summary}

The transition to a 6-layer stackup enables a compact RF front-end design while maintaining proper electromagnetic behavior. Although the 144\,MHz design is relatively tolerant, the 433\,MHz section requires stricter adherence to RF layout practices. Consequently, the layout is designed according to the more stringent UHF requirements to ensure consistent performance across both frequency bands.